\section{自主目标设定与课程学习}
\subsection{目标库与分层 curriculum}
Chelsea 使用“goal replay buffer”记录已达成或接近的目标，结合\textbf{curriculum policy} 选择难度适中的目标进行训练。每次 policy 达成目标后，把该目标与其能力指标一同写入“目标状态 tracker”。

\begin{figure}[H]
\centering
\includegraphics[width=0.88\textwidth]{slides-images/slide-05.jpg}
\caption{Slide 5 讲述如何在 goal-conditioned RL 中实现自动 curriculum。}
\end{figure}

\begin{tabularx}{\textwidth}{lX}
\toprule
阶段 & 核心措施 \\
\midrule
探索 & 采样训练目标时，偏向 novel/uncertain states； \\
平衡 & 设定 success rate window（比如 50\%-80\%）保持学习区； \\
进阶 & 每成功一个目标，使用 parameterized difficulty increment； \\
\bottomrule
\end{tabularx}

\subsection{Goal-Query 与代理思维}
讲者提到“Goal-Query”模块——agent 在训练中不断提问“我应该完成哪个 skill 以便后续推广？”并基于 novelty score 选择目标。这个过程照顾 exploration (novel targets) 与 exploitation (熟悉目标) 的平衡。

\begin{knowledgebox}{Goal-Query 流程}
1. 通过 uncertainty estimate 计算每个 stored goal 的 novelty；\\
2. 如果 novelty 超过 threshold，调度 exploration goal；\\
3. 结合 skill graph 推测 goal 间依赖，决定是否先学习基础 goal。
\end{knowledgebox}

\subsection{Goal space engineering}
Chelsea 强调：goal space 应支持动态缩放，对不同目标设定中心/半径/priority。例如，把 complex goal 拆成 primitives，训练时多次 sample Primitive-level goal，再用 scheduler 提升到 composite goal。

\begin{tabularx}{\textwidth}{lX}
\toprule
策略 & 作用 \\
\midrule
Primitive sampling & 从基础动作组合生成 reward-distinct goals； \\
Composite scheduler & 当 primitive success rate 达标时，按比例加权 composite goal 选择概率； \\
Priority annealing & 根据 novelty/entropy 动态调整 goal priority； \\
\bottomrule
\end{tabularx}

\subsection{本章小结}
自主目标设定把课程学习、Goal-Query 与 difficulty scheduler 结合，使机器人能在没有人类指定任务的情况下逐步积累技能。
\newpage

